\section{Polynomial operators on the primitive cyclotomic expansions}

% BEGIN K82P1
\paragraph{K82P1. Action on all branches and the common constant term.}
Fix either of the two crystalline lattices on one backward arc of
I82B3. Use the transported crystalline basis of J82F2, denote its
space by $D$ over the coefficient field $L$, and write
$F_s^a\subset D$ for the transported filtration at level $s$.
At the primitive $p^{s+1}$ roots use $1+X=\zeta\exp(Y)$.
The coefficient in degree $i\in\mathbb Z$ of a Laurent expansion
is a vector in $B_s\otimes_L D$, where
$B_s=L\otimes_{\mathbb Q_p}\mathbb Q_p(\mu_{p^{s+1}})$.
Over a splitting field these are tuples at all primitive roots.
The physical level module $\mathcal N_{l,s}$ consists of expansions
whose degree-$i$ coefficient lies in $B_s\otimes F_s^{-i}$,
by the filtered Taylor lattice of W79H2--W79H3. We use the
ambient product of tame-invariant expansions, so every component
satisfies the tame action law.

Put $u=1+p$. On the degree-$i$ coefficients the generator acts by
$u^i\sigma$, where $\sigma$ sends the root label $\zeta$ to
$\zeta^u$ by substitution. Likewise a tame element $a\in\mu_{p-1}$
acts by $a^i\tau_a$ with label $\zeta^a$. Indeed the substitution
for either unit $z$ is $\zeta^z\exp(zY)$, and it fixes the
constant crystalline periods. These formulas give equivariance
of the physical expansions and preserve $B_s\otimes F_s^{-i}$.
The automorphisms of the labels commute. The order of $\sigma$
divides $p^s$, since $v_p((1+p)^{p^s}-1)=s+1$. Its powers
and the tame elements generate all units modulo $p^{s+1}$:
$u$ has order exactly $p^s$ there, and the reduction of the
tame elements gives all units modulo $p$. Hence their common
fixed algebra is $L$. Equivalently, the simultaneous fixed
tuples are constant on the transitive set of labels, with
coefficient descent to $L$. The same statement holds with
values in any $L$-vector space, by taking coordinates.

Suppose $(\gamma-1)e$ satisfies the physical level condition.
In the quotient of degree $i$ by $B_s\otimes F_s^{-i}$ the class
of the coefficient is fixed by $u^i\sigma$. Raising this equality
to the $p^s$-th iterate shows that it vanishes for $i\ne0$,
since $u^{ip^s}\ne1$. At zero the class is fixed by both
$\sigma$ and the tame elements. It therefore comes from
$D/F_s^0$. Lift it to $d\in D$. Subtracting the constant $d$
leaves an element of $\mathcal N_{l,s}$. The same $d$ is placed
diagonally in all levels to obtain an element of the prescribed
constant space $D_c$; this membership test concerns just level
$s$. This proves $\mathrm{K2iii}$. The action formulas also
show that $\gamma$ fixes every element of $D_c$ and preserves
$\mathcal N_{l,s}$, proving $\mathrm{DcFix}$.
% END K82P1

% BEGIN K82P2
\paragraph{K82P2. The polynomial for the rounded level module.}
Let $\ell$ denote the length of this arc and $E=\ell-1$, $m=p^E$.
For a finite set of integers $J$ put
\[
 P_J(Z)=\left(\sum_{k=0}^{m-1}Z^k\right)
       \prod_{i\in J\setminus\{0\}}(Z^m-u^{im})\in O[Z].
\]
Negative exponents here are powers of a unit in $O$. At every
level $s<\ell$, K82P1 gives $\sigma^m=1$; thus $\gamma^m$
on the entire degree-$i$ coefficient is $u^{im}$. This identity
holds on all label components simultaneously.

In the notation of A82N1, the rounded coefficient subspace
$X_s(i)$ on a thin letter is zero below zero and $D$ thereafter.
On an up letter it has successive values $0,U_s,D$ with
thresholds $-b_s,0$. On a down letter the thresholds are $0,b_s$
with the same successive values. Here $U_s$ denotes the
rational span of the indicated constant sublattice.
For a Laurent series $f$ the condition
$P_s^g f\in\mathcal N_{b,s}$ means that every coefficient of $f$
lies in $B_s\otimes X_s(i+g)$ in degree $i$. Indeed $P_s$ is
$Y$ times an invertible power series on each branch. Power series
multiplication preserves the coefficient condition by monotonicity
of $X_s$, including multiplication by the inverse of the unit.
Thus $X_s(i+g)$ can exceed $X_s(i)$ only for the $g$ integers
immediately below one of the displayed thresholds.
Take $J_2$ to be the union over the arc of all these intervals.
This is the source's $\mathrm{jumpWindows}$. No interval contains
zero, since $b_s>g$ on thick letters.

If $e=d+f$ with $d\in D_c$ and $P_s^g f\in\mathcal N_{b,s}$,
then $P_{J_2}(\gamma)f$ is zero in each potentially offending
degree, by the corresponding factor and the identity above.
In all other degrees it preserves the required subspace, since
the generator acts only on the label and by a scalar.
Finally $P_{J_2}(\gamma)d=P_{J_2}(1)d$ is constant.
This gives $\mathrm{K3a2}$ for the actual rounded target module.
% END K82P2

% BEGIN K82P3
\paragraph{K82P3. The polynomial for the true filtration.}
For the physical source module take $X_s(i)=F_s^{-i}$.
The same unit argument shows that $P_s^g f\in\mathcal N_{l,s}$
implies $f_i\in B_s\otimes F_s^{-i-g}$. A defect relative
to $F_s^{-i}$ therefore requires a filtration jump at $a$ and
$i=-a-z$ for some $1\le z\le g$. In the jump bounds of I82B3
write $G',S',\eta'$ for the thin, width, and near-zero bounds.
On a thin letter all candidate degrees belong to
$[-G'-g,G')$. On an up or down letter the common near-zero
part is $[-\eta'-g,\eta')$; the other part is respectively
$[-b_s-S'-g,-b_s)$ or $[b_s-g,b_s+S')$. All intervals
denote sets of integers. Their union over the arc is exactly the
source's $J_4=\mathrm{wideWindows}$.

For $e=d+f$ with $d\in D_c$ and $P_s^g f\in\mathcal N_{l,s}$,
the nonzero defect degrees are killed by $P_{J_4}(\gamma)$ as
in K82P2. At zero the second product in that polynomial acts
by $\prod_{i\in J_4\setminus\{0\}}(1-u^{im})$ because
$\sigma^m=1$. The first factor sends $f_0$ to
$\sum_{k=0}^{m-1}\sigma^kf_0$. This is fixed by $\sigma$,
and also by every tame element since the actions commute and
$f_0$ is tame-invariant. By K82P1 it is one common vector in
$D$, including when $L$ does not already split the roots.
Its scalar multiple can be subtracted as a new diagonal
constant. The constant from $d$ is $P_{J_4}(1)d$.
All other coefficients satisfy their required subspaces.
This proves the native $\mathrm{K3a4}$, including a possible
defect at degree zero. No division by the number of branches
is involved in the integral polynomial operator.
% END K82P3

% BEGIN K82P4
\paragraph{K82P4. The integral scalar of the two polynomials.}
For an integer $n\ne0$, the scalar $1-u^n$ is a unit times
$p^{1+v_p(n)}$ already in $\mathbb Z_p$. For a positive integer
prime to $p$ this follows by expanding $(1+p)^n$ modulo $p^2$.
If $v_p(x)\ge1$, the binomial expansion of $(1+x)^p-1$ has
first term $px$ and all later terms of strictly larger valuation,
since $p$ is odd. Iteration, and then inversion of a unit for
negative $n$, proves the formula. Thus
\[
 P_J(1)=\varepsilon_J p^{k_J},\qquad
 k_J=E+\sum_{i\in J\setminus\{0\}}(E+1+v_p(i)),
 \qquad \varepsilon_J\in\mathbb Z_p^\times .
\]
This also handles $E=0$ and empty windows.

Both $k_{J_2}$ and $k_{J_4}$ are at most the original integer
$\kappa=\mathrm{kapPar}(p,g,h,\ell_{\max})$. For $J_2$ use the
common interval $[-g,0)$ and one interval of length $g$ at
each thick letter. For $J_4$ put $c_0=\max(G',\eta')$ and use
the common interval $[-c_0-g,c_0)$ together with the respective
outer intervals in K82P3. There are at most two thick letters,
by I82B3. The common interval has length at most $3g$ and
each outer interval has length at most $S'+g\le2g$. All
these intervals lie in $[-2h,2h]$ since $g\le h$ and
$b_s+S'\le h$. Use empty intervals for missing marks.
They are admissible triples in the exact supremum defining
$\kappa$, with $E=\ell-1\le\ell_{\max}$ and nonnegative
summands. Thus the two native polynomials have product of
evaluations at one equal to $\varepsilon p^k$, where
$k=k_{J_2}+k_{J_4}\le2\kappa$ and $\varepsilon\in O^\times$
is the image of the product of the two units. This is
$\mathrm{PPval}$ in the tested coefficient ring, including
ramified extensions.
% END K82P4
